\documentclass{article}
\usepackage[T2A]{fontenc}
\usepackage[utf8]{inputenc}
\usepackage[russian]{babel}

\begin{document}
\section{Русская типографика и шрифты T2A}

Привет, мир! Это проверка поддержки кириллицы и шрифтов стандарта T2A в Ratex.

Все люди рождаются свободными и равными в своем достоинстве и правах.
Они наделены разумом и совестью и должны поступать в отношении друг друга в духе братства.

\subsection{Размеры шрифтов и метрическое замыкание}

{\footnotesize Сносочный размер (footnotesize): Все люди рождаются свободными и равными.}

{\small Малый размер (small): Все люди рождаются свободными и равными.}

{\normalsize Обычный размер (normalsize): Все люди рождаются свободными и равными в своем достоинстве и правах.}

{\large Увеличенный размер (large): Все люди рождаются свободными и равными в своем достоинстве.}

{\Large Большой размер (Large): Все люди рождаются свободными и равными.}

{\huge Огромный размер (huge): Все люди рождаются свободными.}

\subsection{Начертания в различных размерах}

\normalsize
\textbf{Жирный шрифт: Съешь же ещё этих мягких французских булок, да выпей чаю.}

\textit{Курсивный шрифт: Съешь же ещё этих мягких французских булок, да выпей чаю.}

\textbf{\textit{Жирный курсивный шрифт: Съешь же ещё этих мягких французских булок.}}

\large
\textbf{Жирный увеличенный: Русская типографика T2A.}

\textit{Курсивный увеличенный: Русская типографика T2A.}

\normalsize
\textsf{Рубленый шрифт (Sans): Заголовок или выноска на русском языке.}

\texttt{Моноширинный шрифт: def привет(имя): print(f"Привет, \{имя\}!")}

\end{document}
